\documentclass[12pt]{article}
\usepackage[utf8]{inputenc}
\usepackage[T1]{fontenc}
\usepackage{hyperref}
\title{Yinwum Language Resources for K-12 Curriculum Development}
\author{Compiled by Claude Sonnet 5.5 (claude-sonnet-5-5) for Andr\'as Kornai}
\date{}
\begin{document}
\maketitle

\section{Introduction}
Yinwum (ISO 639-3 yxm; AIATSIS code Y29) is an extinct Paman language of Cape York Peninsula, Queensland. Its people lived around the upper Wenlock River. Search results say it was no longer spoken by the 1960s. The sources found are encyclopedic and describe the language only briefly. No teaching materials were found, and what follows is a short report, not a curriculum bibliography.

\section{Primary Educational Resources}
No primers, textbooks, school programmes or classroom materials for Yinwum turned up in search results. A heritage or revitalization curriculum would have to be built from archival sources and, ideally, with the descendant community.

\section{Linguistic and Reference Materials}
\begin{itemize}
\item Wikipedia, ``Yinwum'': a short overview covering location, extinction, phonology and neighbouring groups. It is a secondary source and should be checked against primary ones. \url{https://en.wikipedia.org/wiki/Yinwum}
\item Hgcosmos PDF on Yinwum: a short cultural note surfaced in search results. I did not verify its contents. \url{https://www.hgcosmos.org/assets/culture/AZ/Y/Yinwum.pdf}
\end{itemize}
The search summary names Ursula McConnel's 1939 work on the social organization of Cape York tribes as the main ethnographic source. No link to it was found.

\section{Religious and Community Resources}
No scripture translations or community organisations specific to Yinwum were found.

\section{Digital and Audio Resources}
No audio, video or corpus resources were found in search results.

\section{Curriculum Development Context}
Little documentation exists, and the search found no surviving speakers or existing programme. Any work should begin by consulting AIATSIS and the relevant Cape York Aboriginal communities. It should also check archival wordlists and avoid inventing vocabulary.
\end{document}